\subsection{Step-size and noise sensitivity}
\label{sec:dt-noise}

Two more one-factor sweeps vary the data rather than the model
(Figure~\ref{fig:noise-dt}, Table~\ref{tab:noise-dt}). The \textbf{step-size
sweep} changes the sampling interval $\Delta t$, and with it the number of
training points and the solver step $h=\Delta t/2$. The trend is not
monotone. $\Delta t=0.1$ extrapolates best; $\Delta t=0.05$ is $1.5\times$
worse with twice the data and twice the cost; and $\Delta t=0.2$ reaches the
\emph{lowest} training loss and the worst extrapolation ($5.1\times$). The
work-precision data rule out truncation as the cause: Tsit5's own error on
the true field at these three steps is $2.6\times10^{-18}$, $4.7\times10^{-14}$
and $2.2\times10^{-10}$ in MSE, at least six orders of magnitude below any
learned error. The coarse grid hurts extrapolation because 19 samples
constrain the field less (it fits them more closely and generalises worse).
The fine grid hurts it because the same 10{,}000 epochs buy a similar
training loss on a harder, longer unrolled graph. Here, step size is a
data-density choice, not a numerical-accuracy one.

\begin{table*}[t]
\centering
\footnotesize
\renewcommand{\arraystretch}{1.1}
\caption{Step-size and observation-noise sweeps (Tsit5, RBF, $10^4$ epochs).
Noise is added only to the training window; all errors are measured against
the clean truth.}
\label{tab:noise-dt}
\begin{tabular}{@{}L{3.4cm} C{1.8cm} C{2.3cm} C{2.3cm} C{1.9cm} C{1.4cm} C{2.0cm}@{}}
\toprule
\hdrrow \thh{Configuration} & \thh{Training points} & \thh{Training MSE} & \thh{Extrap.\ MSE} & \thh{Extrap.\ $R^2$} & \thh{$\hat L$} & \thh{Wall-clock}\\
\midrule
$\Delta t=0.05$ ($h=0.025$) & 71 & $9.37\times10^{-5}$ & $1.34\times10^{-4}$ & $0.99995$ & $7.46$ & 3425 s\\
$\Delta t=0.1$ (reference) & 36 & $8.83\times10^{-5}$ & $\mathbf{8.92\times10^{-5}}$ & $0.99997$ & $6.91$ & 1713 s\\
$\Delta t=0.2$ ($h=0.1$) & 19 & $7.83\times10^{-5}$ & $4.51\times10^{-4}$ & $0.99985$ & $6.64$ & 882 s\\
\midrule
$\sigma=0.01$ & 36 & $1.23\times10^{-4}$ & $1.22\times10^{-3}$ & $0.99958$ & $7.18$ & 1686 s\\
$\sigma=0.05$ & 36 & $7.41\times10^{-4}$ & $2.54\times10^{-2}$ & $0.99129$ & $8.69$ & 1672 s\\
$\sigma=0.1$ & 36 & $2.28\times10^{-3}$ & $1.37\times10^{-1}$ & $0.95282$ & $9.55$ & 1534 s\\
\bottomrule
\end{tabular}
\end{table*}

\begin{figure*}[t]
\centering
\begin{minipage}[t]{0.45\textwidth}
\centering
\includegraphics[width=\linewidth]{figures/ablation/02_error_vs_stepsize_loglog.png}
\end{minipage}\hfill
\begin{minipage}[t]{0.45\textwidth}
\centering
\includegraphics[width=\linewidth]{figures/analysis/noise_sensitivity.pdf}
\end{minipage}
\caption{Left: error against sampling interval, with $O(\Delta t)$ and
$O(\Delta t^2)$ reference slopes; changing $\Delta t$ also changes the
number of training samples. Right: error against observation-noise level,
with least-squares power-law fits; the dotted line is the noise variance
$\sigma^2$.}
\label{fig:noise-dt}
\end{figure*}

The \textbf{noise sweep} adds Gaussian noise of standard deviation $\sigma$
to the 36 training observations. Across $\sigma\in\{0.01,0.05,0.1\}$ the
extrapolation MSE follows a clean power law,
$\mathrm{MSE}_{\text{extrap}}\approx12.9\,\sigma^{2.02}$. In other words,
extrapolation error is proportional to the noise variance, with a gain of
about $13$. The training-window error against the clean truth grows more
slowly ($\propto\sigma^{1.24}$) and sits below $\sigma^2$ at every level
($0.23\sigma^2$ at $\sigma=0.1$), so the ODE constraint does denoise the
fit inside the window. What it cannot do is stop the remaining field error
from compounding into phase error beyond the window. The field also
roughens as it absorbs the noise ($\hat L$ rises from $6.91$ to $9.55$).
Section~\ref{sec:track-e} compares this sensitivity against sparse symbolic
regression.
